Apply the categorical hidden-functional sharp-region theorem with
\(\mathcal Z=\mathcal B=\mathcal X=\mathcal Y=\{0,1\}\), with the four
vertices and edges specified by \(G_{\mathrm w}\), observed set
\(V_{\mathrm w}\), functional set \(W_{\mathrm w}\), singleton treatment
set \(T_X\), and outcome set \(Y_{\mathrm w}^{\star}\).  The law
\(p_{\mathrm w}\) is feasible because it is, by definition, induced by a
model in \(\mathfrak M_{\mathrm w,X}\), and the assumed inequalities
\(p_{\mathrm w}(X=x)>0\) are precisely singleton joint-treatment
positivity for the two treatment values.

For binary \(Y\), a probability vector is determined by its coordinate at
\(Y=1\), and
\[
 \{L_x+M_xs:s\in\Delta(\{0,1\})\}_{Y=1}
 =[L_x(1),L_x(1)+M_x].
\tag{1}
\]
The categorical lower coordinate is
\[
 L_x(1)
 =\sum_{z:p_{\mathrm w}(z,x)>0}
 p_{\mathrm w}(z)\frac{p_{\mathrm w}(z,x,Y=1)}{p_{\mathrm w}(z,x)},
\tag{2}
\]
where every division is guarded by \(p_{\mathrm w}(z,x)>0\), and its
categorical missing mass is
\[
 M_x=\sum_{z:p_{\mathrm w}(z)>0,\ p_{\mathrm w}(z,x)=0}
 p_{\mathrm w}(z).
\tag{3}
\]
Equations (2)--(3) are exactly the quantities in
Definition~\(\mathrm{def:binary\text{-}identified\text{-}region}\).
Taking the Cartesian product of (1) for \(x=0,1\) therefore gives
\[
 \left\{\bigl(Q_{X,M}(0,1),Q_{X,M}(1,1)\bigr):
 M\in\mathfrak M_{\mathrm w,X},\ p_M=p_{\mathrm w}\right\}
 =[L_0,L_0+M_0]\times[L_1,L_1+M_1]
 =\mathcal I(p_{\mathrm w}).
\tag{4}
\]
The simultaneous-attainment part of the categorical theorem uses one
common slack distribution for all functional-profile rows missing a given
\(x\), independently across \(x\), and never changes an alphabet.  Hence
every point of the rectangle in (4) is attained by a model on the original
four binary alphabets.

Because \(\lvert\mathcal Y\rvert=2\), the categorical collapse criterion
says that (4) is a singleton if and only if \(M_0=M_1=0\).  Every summand
in (3) is strictly positive, so \(M_x=0\) is equivalent to
\(p_{\mathrm w}(z,x)>0\) for every \(z\) with
\(p_{\mathrm w}(z)>0\).  Requiring this for both \(x=0,1\) proves the
stated point-identification equivalence.